\begin{center}
    {\large \textbf{Giorno 11} \textbar\ La città dei ciapa puer \textbar\ 19 Agosto 2024 \textbar\ \textbf{Ubud}, Bali}
\end{center}

\noindent\rule{\textwidth}{0.4pt}
\begin{figure}[h!] % Portrait images below
    \centering
    % First portrait image (on the left)
    \begin{minipage}[h]{0.48\textwidth} % Set width equal to half the page
        \centering
        \includegraphics[width=\textwidth]{images/flags/Screenshot Ubud.png}
    \end{minipage}
    \hfill
    % Text
    \begin{minipage}[h]{0.48\textwidth}
        \raggedright
        \textbf{Superficie}: 4.567 km² \\
        \vspace{0.3cm}
        \textbf{Popolazione}: 3,3 milioni \\
        \vspace{0.3cm}
        \textit{L'isola}
    \end{minipage}
\end{figure}

\landscapeLargeMargin{images/day11/PXL_20240819_010056305.jpg}

\landscapeLargeMargin{images/day11/PXL_20240819_021835421.jpg}

\landscapeLargeMargin{images/day11/PXL_20240819_022517989.jpg}

\landscapeLargeMargin{images/day11/PXL_20240819_092049939.jpg}

\landscapeLargeMargin{images/day11/PXL_20240819_092104408.jpg}

\landscapeLargeMargin{images/day11/PXL_20240819_091836542.jpg}

\landscapeLargeMargin{images/day11/PXL_20240819_094406795.jpg}

\landscapeLargeMargin{images/day11/PXL_20240819_104207087.jpg}

\landscapeLargeMargin{images/day11/PXL_20240819_104210712.jpg}

\landscapeLargeMargin{images/day11/PXL_20240819_111157635.jpg}

\landscapeLargeMargin{images/day11/PXL_20240819_111242103.MP.jpg}

\landscapeLargeMargin{images/day11/PXL_20240819_113145698.jpg}

\landscapeLargeMargin{images/day11/PXL_20240819_112342862.NIGHT.jpg}

\noindent 19 Agosto 2024 \newline

Notte di riposo sereno: credo di essermi finalmente abituato al fuso orario locale. Oggi mi servirà essere riposato perché ci attende un lungo viaggio.

Questa mattina optiamo per la colazione al ristorante dell'hotel, così da esplorare anche quell'ambiente — il servizio in camera, dopo un inizio degno di nota, ha lasciato abbastanza a desiderare. Il ristorante italiano all'interno del Villas Edenia, L'Osteria, ricorda molto il ristorante del Pearl Beach sulla Gili Asahan: completamente all'aperto, con mura di legno chiaro, decorazioni di corda e molto bambù. È situato su una piattaforma di legno a cui si accede tramite tre scalini dopo essersi tolti le scarpe, e dispone di un'area relax con divanetto e una piccola libreria popolata dai libri lasciati dai viaggiatori.

Ci sediamo e ordiniamo: il solito wrap per Elisa e una colazione diversa per me, l'Edenia Pancake, un pancake alla banana ricoperto di gelato al fiordilatte e guarnito con anguria e frutto della passione. Ottimo, ma un po' povero.

Terminata la colazione torniamo in camera e mentre Eli finisce di prepararsi inizio a portare i miei bagagli alla reception per il check-out. Qui lo staff del resort completa il crescendo di incompetenza di questi tre giorni con una ciliegina sulla torta.

Prima mi chiede il numero preciso di birre consumate dal frigo-bar — come se tenere il conto di quanta birra si beve fosse una cosa facile! — poi mi dice che non posso pagare con la carta per un motivo incomprensibile, che potenzialmente è pigrizia, dal momento che mi sta facendo fattura di tutto. Pago controvoglia rimanendo praticamente senza contanti, con appena il necessario per pagare il calesse.

Per quest'ultimo sono fiducioso: l'ho prenotato il giorno prima e stamattina, appena arrivato alla reception, ho ricordato della prenotazione al ragazzo che ha annuito in risposta. Esattamente alle 9:20, spaccando il minuto, Elisa arriva con i suoi bagagli. Del cavallo, neanche l'ombra.

Chiedo al ragazzo alla reception che chiama qualcuno col walkie-talkie e mi rassicura. Aspettiamo. Aspettiamo. Aspettiamo. Dopo dieci minuti vado a insistere ricordando che dobbiamo prendere il traghetto e non possiamo arrivare tardi. Riesco finalmente a contagiarlo con la mia tensione, parla concitato con qualcuno e mi dice che il calesse è in arrivo. Intanto i minuti passano.

Elisa insiste per andare a piedi e intercettare un passaggio lungo la strada, mentre io cerco di tenere la posizione. Sono quasi le 9:40 e dico al ragazzo che se entro quell'ora il nostro passaggio non fosse arrivato ce ne saremmo andati. Lui cerca di tranquillizzarmi dicendo che tanto il traghetto non parte alle 10 ma alle 11 — frase che scatena in me pensieri di brutalità e sevizie. Prima di tutto, non me ne frega niente se ci sono altri traghetti che partono dopo: potrei avere un appuntamento a destinazione che non mi consente di perdere un'ora ad aspettare un cavallo. Secondo, e più importante, ormai non mi fido più di lui.

Alle 9:40 spaccate prendo Elisa e ce ne andiamo senza nemmeno salutare. Lo sento che intanto ha chiamato nuovamente qualcuno al walkie-talkie e, ormai in lontananza, ci prega di aspettare. Anche Elisa, che per prima insisteva per andare via, ci rimane un po' male per non aver salutato — ma io ormai sono nero, e vado via senza nessun rimorso.

Percorriamo a passo svelto il tratto fino allo stradone principale, carichi con borsoni e zaini, e qui troviamo subito un calesse che in un attimo ci carica e ci porta a destinazione. Durante il tragitto incrociamo ben quattro carrozze vuote, di gran fretta, che sono quasi certo fossero quelle chiamate dal ragazzo alla reception del Villas Edenia. Spero non riceva una strigliata esagerata per essere stato troppo superficiale.

Arriviamo alla biglietteria pochi minuti prima della partenza della barca, mostriamo la prenotazione e riceviamo due tesserini come biglietti. Corriamo al porto dove, grazie a Elisa, superiamo tutti per imbarcarci sulla prima barca pronta a salpare. All'ultimo checkpoint però siamo fermati da un guardiano che ci comunica che dobbiamo pagare una tassa — forse quella per entrare a Bali — e che dobbiamo quindi passare dagli uffici della capitaneria. Gli diciamo che non abbiamo tempo e chiediamo se possiamo pagare lui, che acconsente ma ovviamente "just cash." Dopo aver pagato resort e calesse sono senza contanti e scopriamo che anche gli uffici non accettano le carte, quindi dovremmo andare a prelevare. Non esiste, così insistiamo e il guardiano accetta un pagamento in euro, per il quale — stranamente — non ha il resto. Va bene così: paghiamo le tasse per cinque persone e salviamo questa voce di costo come "mazzetta."

All'ultimo controllo ci chiedono di lasciare i bagagli da mettere nella stiva insieme a quelli di altre centinaia di persone: Elisa sgattaiola in avanti con tutte le borse addosso ignorando i richiami, mentre io faccio il diplomatico e alla richiesta del perché non voglia lasciare il mio bagaglio rispondo che contiene "medicines and electronics." Mi lasciano passare senza insistere, consapevoli che tanto a pagarla sarei stato io.

Salgo a bordo chinandomi per accedere attraverso il boccaporto e per poco non volo in avanti per tutto il peso che ho addosso. Ricorso che sto ancora portando lo scatolone con i vasi di terracotta. Ci sediamo vicino a una coppia di italiani che ha lasciato un paio di posti liberi perché bagnati. Mi sembra una buona idea per stare vicino a Eli, così metto un asciugamano sul bagnato e mi siedo. Siamo quasi gli ultimi a salire e salpiamo subito.

Da quello che so il viaggio dovrebbe durare un'ora e mezza, anche se sul biglietto sono indicate tre ore. Dopo i primi dieci minuti realizzo che sarà il viaggio della speranza: sono già salito sudato, nel mio posto si muore dal caldo e non circola un filo d'aria, il finestrino perde e l'acqua delle onde mi gocciola in testa e sui vestiti. Inizio a sudare copiosamente ed Elisa si preoccupa. Cerco di fare un po' di training autogeno per rilassarmi e abbassare i battiti, ma i miei sforzi sono vani ed Elisa è sempre più preoccupata, col risultato che io inizio a minimizzare: "sarà solo un po' di mal di mare unito al caldo."

Passa quasi un'ora e la situazione non accenna a migliorare, anzi. Elisa mi sfrega sulla fronte e sul collo delle salviette umide per cercare di rinfrescarmi. Le impedisco di rovesciarmi acqua sulla testa perché ho già metà del corpo fradicio per via del finestrino, cosa che non fa altro che aumentare il mio nervoso goccia dopo goccia.

Dopo oltre un'ora di sofferenza — a detta di Eli sono diventato bianco e le mie labbra sono diventate blu — decido finalmente di ascoltarla e mi sposto sui gradini davanti al portellone del boccaporto, una fila più indietro, dove l'afflusso d'aria è maggiore. Una coppia di pensionati australiani mi osserva sconvolta e io incolpo il finestrino che perde, mostrando la manica e la gamba del pantalone bagnate. Loro ridono, sostenendo che pensavano fosse tutto sudore. Gli rispondo “fifty-fifty”.

La temperatura corporea inizia a scendere gradualmente, ma al contempo sale più forte la nausea da mal di mare. Sento che sto per vomitare e chiedo a Elisa un sacchetto. È passata ormai un'ora e mezza, è chiaro che il viaggio durerà tre ore e che siamo solo a metà strada, e sono sicuro di non riuscire a resistere fino alla fine. Mi viene in mente quando, durante una lezione di apnea in cui avevo patito il mal di mare, gli istruttori mi avevano consigliato di fissare un punto fermo a riva senza distogliere lo sguardo. Ci provo: vedo da lontano la punta del Monte Agung che si staglia all'orizzonte tra le nuvole sopra Bali e mi concentro con tutto me stesso per non levare lo sguardo neanche per un secondo. Elisa nel frattempo continua a curarmi a suon di biscotti e Biochetasi. L'insieme delle cose fa sì che mi senta un po' meglio.

L'ultima ora di viaggio prosegue in questo modo, con me immobile a scrutare l'orizzonte. Dopo tre interminabili ore arriviamo a destinazione, con una breve sosta di fronte al porto in attesa del nostro turno di attraccare. Ringrazio di non aver lasciato il bagaglio: anche se mi ha quasi fatto volare all'ingresso e prendere una testata all'uscita, questa scelta ci ha permesso di guadagnare probabilmente un'ora buona, visto il tempo necessario per scaricare i bagagli di tutti i passeggeri.

Cerchiamo di uscire dal porto ma il caos regna sovrano e ci mettiamo cinque minuti buoni a percorrere venti metri, tra la fila di chi parte e quella di chi arriva. Entrati nella cittadina portuale di Padang Bai scopriamo che questa è ancora più congestionata del porto. Centinaia di auto sono in coda e da tutte le parti sentiamo gridare "Taxi! Taxi!" in maniera decisamente aggressiva, così che mentre cerchiamo di allontanarci dalla calca veniamo inseguiti più volte da uno stesso tassista che vuole a tutti i costi farsi ingaggiare.

Sono ormai esausto: gli dico che degli amici mi devono venire a prendere, al che lui si mette a ridere dicendo di non crederci. Sento di aver rischiato la vita sulla barca, sono disidratato, ho mal di testa e sono carico come un mulo — le risate di questo invadente sconosciuto sono la goccia che fa traboccare il vaso. Faccio capire al tassista, con tono educato ma condito da una certa dose di volgarità, e con una calma solo apparente mentre gli occhi gli urlano addosso tempesta, che da lui non avremmo mai accettato un passaggio e che era ora che se ne andasse. Stavolta afferra il messaggio.

In teoria le app di trasporto privato Grab e Gojek qui non sono ammesse, ma riusciamo a rimediare un passaggio su Gojek, che secondo Enrico è il più utilizzato nella zona di Ubud. Troviamo subito un driver che ci invita a raggiungerlo a inizio paese, poiché se avesse tentato di venire da noi avrebbe rischiato di rimanere bloccato nel traffico. Accettiamo anche quest'altro scherzo del destino e ci prepariamo a camminare oltre un chilometro a piedi con borse e zaini, in salita e sotto il sole dell'una. La mia camicia ormai è da buttare, da ogni macchina continuano a gridare "Taxi! Taxi!", la strada è spesso interrotta da enormi buche e — come se non bastasse — passiamo sotto una cava dalla quale piovono massi che per fortuna si fermano prima di arrivare sulla strada.

Finalmente raggiungiamo il nostro autista, che se ne sta coricato nel bagagliaio dell'auto. Carichiamo, saliamo e la voglia di fare due chiacchiere con lui è del tutto assente, così rispondo alle domande di rito senza mostrare intenzione a proseguire nella conversazione.

Il viaggio è tranquillo, con la solita guida a base di clacson e sorpassi azzardati a cui ormai siamo abituati. Tengo sotto controllo la strada con il telefono — non so se per una questione di controllo o per evitare che ci rapisca a nostra insaputa — quando a un certo punto mi rendo conto che le scorciatoie che ha preso ci stanno portando in tutt'altra direzione. Realizzo di aver selezionato la destinazione sbagliata: un hotel con nome analogo al nostro ma dall'altra parte di Ubud. Interrompo l'autista e gli chiedo di cambiare strada, promettendogli di pagare la differenza, senza sapere che avrei potuto cambiare la destinazione direttamente dall'app. Lui è molto gentile e comprensivo e mi chiede solo di pagargli la differenza in contanti, ma quando gli chiedo una cifra risponde "up to you." Ah, la gentilezza indonesiana.

Arriviamo a Ubud e cerco subito un ATM, dal momento che sono ancora senza contanti. Il primo in cui ci fermiamo ovviamente non funziona, così ne cerco un altro e non trovandone uno nelle vicinanze dell'hotel mi faccio lasciare a un money exchange e cambio lo stretto necessario per pagarlo. Quando finalmente scendiamo dall'auto gli do una banconota da 100.000 rupie, circa 5,80€ — praticamente la metà della corsa di novanta minuti che ha fatto per portarci fin lì. Quando vede la banconota fa un'espressione incredula e intuisco di aver esagerato. Amen. Non connetto più, voglio solo un letto.

\begin{center}
    % Decorative line
    \noindent\rule{0.6\textwidth}{0.4pt}
\end{center}

Arriviamo finalmente al nostro nuovo alloggio, il Gayatri. Chiamarlo "albergo" sarebbe riduttivo: è come varcare la soglia di un giardino dell'Eden. Siamo accolti da donne induiste vestite con abiti di pizzo bianco iridescente, unghie scintillanti e il tradizionale puntino rosso sulla fronte. Fin da subito ci rendiamo conto di essere entrati in un mondo che segue regole diverse: scompaiono veli e tuniche monacali, lasciando spazio a colori vivaci e capelli raccolti con eleganza, ornati da fiori di loto rosa. Sono adornate con collane, bracciali, anelli, sandali brillanti e ciglia finte, eppure mantengono intatta la loro identità culturale. Gli abiti, pur ricchi e sfarzosi, restano profondamente diversi da quelli occidentali. I sorrisi sono ovunque e la cordialità è disarmante. Gli uomini indossano Sarong color senape e turbanti bianchi abbinati a camicie tono su tono. Tutto è incredibilmente ordinato.

La nostra villa è una meraviglia: un patio sorretto da colonne, tetto in bambù, un lampadario orientale dorato in ferro battuto e ceramica dipinta a mano. Le modanature decorano ogni angolo, mentre il pavimento è interamente in marmo bianco. Le pareti in pietra scolpita sono un tripudio di arte: draghi alati dalle code interminabili, foglie di giglio, scimmie, serpenti, palme, volti rituali e fiori scolpiti con una cura tale da sembrare vivi. Tutto evoca davvero l'immagine di un paradiso perduto. Le finestre e le porte sono incorniciate da intagli raffinati, con laccature bordeaux e fiori di loto dorati. La porta della camera è in legno naturale, intagliata con eleganza e rifinita da un bordo rosso iridescente. Le maniglie in ottone e il tavolino con sedie lavorate a mano completano un'atmosfera incantata. Davanti a noi, un laghetto punteggiato di ninfee e popolato da statue di draghi e divinità completa questo scenario da sogno.

Per questa descrizione ringrazio ovviamente Elisa — immagino si sia capito che non era opera mia, dal momento che neanche so cosa sia una modanatura.

Sono tutti estremamente gentili e dopo il check-in ci accompagnano in stanza dove un saluto di fiori ci attende oltre l'ingresso: "HI :)". La stanza è spaziosa e molto luminosa, ma quello che mi serve in questo momento è solo un letto su cui buttarmi.

Purtroppo non c'è tempo per dormire: tutto il nostro tempo a Bali è ancora da organizzare e, dal momento che la visita a Kawah Ijen è saltata perché quest'anno il vulcano è in attività, abbiamo anche una notte libera da prenotare. Mi occupo prima del problema principale e dopo un breve giro su Booking trovo un bungalow libero esattamente a fianco della nostra struttura attuale, a un decimo del costo grazie a uno sconto dell'ultimo minuto. Probabilmente qualcuno ha cancellato e si sono ritrovati in alta stagione con una camera da riempire. Sistemato. Dobbiamo poi organizzare un tour per l'indomani, così dopo una veloce doccia ci addentriamo per Ubud alla ricerca di un tour operator. Questo piano fallisce miseramente, perché perderò completamente il controllo di mia moglie.

Elisa a Ubud va fuori di testa. Questo posto per lei è il Paese dei Balocchi: negozi di arte e scultura indonesiana, caffè e ristoranti all'aperto con atmosfera chill, decorazioni di bambù dappertutto, bancarelle con i souvenir più bizzarri. Non riesco a tenerla ferma — appena vede qualcosa che la attira ci si fionda alla velocità della luce. Soprattutto non riesco più a farmi ascoltare: le dico una cosa, annuisce e poi fa tutt'altro, è come se mi si fosse rotta. Per me questo posto è il più grande mercato di ciapa puer al mondo. La quantità immonda di cagate in vendita è inverosimile e so già che sarà una dura lotta per non tornare a casa con le valigie piene di statue e posate. Elisa inizia a scattare fotografie a raffica — mi viene in mente una citazione da un libro di Irvine Welsh: "come un AK-47 nelle mani di un epilettico."

Cerco informazioni sui tour e chiedo supporto a Eli, che però vuole andare a vedere le statue di Ganesh giganti. Le chiedo se prima di andare al localino con la terrazza piena di palme può accompagnarmi a chiedere dei tour dei templi — mi dice di sì e non appena mi giro è già davanti al localino a scattare foto e a sbirciare il menù.

Mi arrendo. Sono anche troppo stanco per affrontare la situazione oggi, così decido che l'indomani l'avremmo trascorso ancora a Ubud: visitare qualche tempio nei paesi limitrofi, vedere la Foresta delle scimmie e organizzare con calma i tour dei templi più lontani, una volta sedata mia moglie.

A proposito di scimmie: mentre percorriamo Monkey Forest Road in direzione della foresta vediamo già molte scimmie che si aggirano sui tetti delle case e sui cavi dei tralicci. Dopo aver comprato degli snack per rifocillarci, assistiamo alla scena di un cane che cerca di mordere la coda di una scimmia, la quale si rifugia sul tetto di un cafè. La proprietaria si gode divertita la scena e mi spiega che, se non fosse stato per il cane, la scimmia mi avrebbe sicuramente già rubato la scatola di biscotti che tenevo in mano. Divertito, lo racconto subito a Elisa intenta a scattare primi piani alla scimmia. Sembra ascoltarmi, ma quando arriviamo nei pressi della foresta veniamo attaccati da una scimmia che vuole il sacchetto di patatine che ho in mano. Devo buttarle tutte e una mezza dozzina di scimmie si avventa sul bottino mentre io scappo. Dopo aver osservato che effettivamente la signora del caffè ci aveva avvisato, Elisa mi risponde: "Quale signora?". Non ce la posso fare.

Camminiamo ancora, passando davanti alla foresta al centro di Ubud e procedendo oltre a passo di lumaca, scattando foto a ogni piastrella e ogni sasso. A un certo punto si fa buio ed entriamo in un supermercato che dovrebbe avere un bancomat — perché in tutto questo siamo ancora senza rupie. Finalmente trovo l'ATM e prelevo, ma quando raggiungo Elisa questa ha già riempito un cestino di prodotti da comprare e portare a casa, principalmente tè e caffè. Sono vicino all'esaurimento, anche perché a parte lo spuntino di biscotti e patatine sono a stomaco vuoto da stamattina. Convinco Eli a posare tutto — consapevole che me l'avrebbe fatta pagare — e usciamo per dirigerci al Caffè Wayan, un locale storico che le ha consigliato la sua amica Lory.

Breve storia del Caffè Wayan, estrapolata direttamente dalla prima pagina del loro menù: la storia di questo locale risale al 1980, quando Ibu Wayan aveva un piccolo warung, una bancarella di cibo in bambù. Ibu Wayan vendeva caffè ai contadini e preparava un budino di riso nero sulla strada che oggi è la Monkey Forest Road. "Ibu" significa "madre" in indonesiano, e la sua passione per la cucina è sempre stata evidente. Con l'aumento del turismo, nel 1986 ha aperto ufficialmente il Café Wayan, con sei tavoli e due semplici bungalow per gli ospiti. Ibu, desiderosa di imparare ricette da diverse culture, ha ampliato il menu per includere piatti occidentali per soddisfare anche i visitatori da molto lontano. Suo marito Ketut, un artista tradizionale, ha iniziato a progettare giardini e creare alloggi per gli ospiti.

Chiediamo di poterci accomodare in un gazebo anziché ai tavoli interni e una cameriera in abiti tradizionali ci accompagna attraverso un lungo giardino punteggiato di gazebo e strutture eleganti. Ogni angolo è pensato per offrire un'esperienza di cena zen: silenziosa, riservata, immersa nella quiete. Gli ospiti, pur restando all'aperto, sono perfettamente isolati gli uni dagli altri in un'atmosfera di pace e discrezione.

Il nostro gazebo si trova in fondo al giardino, affacciato su un laghetto punteggiato di ninfee galleggianti e circondato da statue scolpite nella pietra lavica. È un piccolo gioiello in perfetto stile balinese, interamente in legno, adornato con decorazioni rosse e dorate, con una luce soffusa che rende l'atmosfera intima e suggestiva. Saliamo sul gazebo e ci sediamo a gambe incrociate su un tappeto di iuta con cuscini, davanti a un tavolo in teak con al centro una candela profumata.

Essendo un posto tradizionale ordiniamo piatti molto tipici. Come antipasti, i Tempe — fagiolini di soia fritti in spezie balinesi — e i Bregedel, crocchette di mais in pastella d'uovo. Come primi, un Mie Goreng e un Nasi Campur, piatto tradizionale a base di riso bianco, pollo al curry, pollo Satay, uova allo zafferano, cocco saltato, tempe e bregedel. Il tutto accompagnato da frullati di frutta per reintegrare un po' di vitamine. Il cibo è delizioso e facciamo quasi fatica a finirlo viste le quantità. Prima di andarcene ordiniamo un altro frullato e una tisana — tra la giornata piena e l'umidità di Bali ci sentiamo molto disidratati. Elisa ci tiene a elencare gli ingredienti della sua tisana: spicchi di curcuma, bastoncino di cannella, punte di anice stellato, lime, zucchero di roccia, foglia di alloro fresca e chiodi di garofano.

Terminata la cena ci concediamo qualche scatto di rito: il luogo è davvero unico, qualcosa che non avevamo mai visto prima. I gazebo isolati, i pavimenti decorati a mosaico, gli stagni pieni di carpe koi, le piante di bonsai e le palme ovunque creano un'atmosfera surreale.

Poi ci avviamo verso casa, ma una strana luce verde nel cielo cattura la nostra attenzione. Chiediamo a una cameriera di cosa si tratti e con un sorriso ci spiega che sono LED applicati a un aquilone: a Bali il cielo si riempie spesso di aquiloni altissimi, usati per scopi rituali. Un ultimo tocco di magia prima della buonanotte.

Ci incamminiamo verso casa non senza fatica: l'umidità si fa sentire, la stanchezza pesa e i negozi di souvenir ci rallentano con le loro mille luci e colori. Ci lasciamo volentieri attrarre dalla musica live che esce da alcuni locali lungo la strada, fermandoci per qualche minuto ad ascoltare.

Rientrati al Gayatri ci fermiamo ad esplorarlo con la magia delle luci notturne. All'ingresso uno stagno popolato di carpe koi e adornato da statue ci dà il benvenuto. Oltre le sistemazioni degli ospiti si apre un'enorme piscina illuminata da luci rosse, con lettini coperti da teli e grandi ombrelloni con LED accesi fino al mattino. Poco più avanti si trova quella che sembra essere l'area colazione: un grande gazebo per sedersi in stile tradizionale a gambe incrociate, un'area centrale con tavolini rotondi e sgabelli in pietra, e una zona più occidentale con un lungo tavolo in legno circondato da tre immense librerie colme di libri in ogni lingua. Seguendo un sentiero arriviamo in un vasto prato dove si trova una grande vasca in pietra con i segni di un falò acceso in precedenza, circondata da ombrelloni con luci LED blu. Più avanti ancora, una grande capanna in bambù, probabilmente usata come palco per spettacoli o eventi, chiude il percorso.

Esploriamo tutto con curiosità, scattiamo foto, ci affacciamo ovunque possiamo, frenando a fatica la tentazione di addentrarci nel prato con l'erba alta o nella foresta che lo costeggia. Poi, soddisfatti, torniamo in camera.

Ci sediamo un momento sul patio a sorseggiare due grandi bicchieri d'acqua fresca — io, ancora assetato, termino subito la prima bottiglia. Prima di dormire ordiniamo la colazione per il giorno dopo tramite WhatsApp e riceviamo una conferma immediata. Efficientissimi… sperando che non si rivelino come quelli della Gili T.

Ci infiliamo finalmente nel nostro letto a baldacchino, pronti a crollare tra sogni e zanzariere.

\pagestyle{empty} % disable numbers

%coppia
\landscapeLargeMargin{images/day11/PXL_20240819_085232832.jpg}
\landscapeLargeMargin{images/day11/PXL_20240819_094344225.MP.jpg}
%

%coppia
\portraitFullscreen{images/day11/PXL_20240819_112439928.NIGHT.jpg}
\landscapeLargeMargin{images/day11/PXL_20240819_121439071.PORTRAIT.jpg}
%

%coppia
\twoImagesLargeMargins{images/day11/PXL_20240819_125326824.jpg}
                      {images/day11/PXL_20240819_125417409.MP.jpg}
\portraitFullscreen{images/day11/PXL_20240819_130242984.MP.jpg}
%

%coppia
\twoImagesLargeMargins{images/day11/PXL_20240819_125417409.MP.jpg}
                      {images/day11/PXL_20240819_125422118.jpg}
\twoImagesLargeMargins{images/day11/PXL_20240819_125537221.jpg}
                      {images/day11/PXL_20240819_125915956.MP.jpg}
%

%coppia
\landscapeLargeMargin{images/day11/PXL_20240819_131948473.MP.jpg}
\landscapeLargeMargin{images/day11/PXL_20240819_131901004.MP.jpg}
%

%coppia
\splitImageLeft{images/day11/PXL_20240819_132029821.jpg}
\splitImageRight{images/day11/PXL_20240819_132029821.jpg}
%

%coppia
\twoImagesLargeMargins{images/day11/PXL_20240819_132453581.MP.jpg}
                      {images/day11/PXL_20240819_132130861.jpg}
\landscapeThinMarginCentered{images/day11/PXL_20240819_150229658.jpg}
%

%coppia
\twoImagesLargeMargins{images/day11/PXL_20240819_150339668.MP.jpg}
                      {images/day11/PXL_20240819_150632060.jpg}
\twoImagesLargeMargins{images/day11/PXL_20240819_150944555.NIGHT.jpg}
                      {images/day11/PXL_20240819_151236329.NIGHT.jpg}
%

%coppia
\splitImageLeft{images/day11/PXL_20240819_151811404.MP.jpg}
\splitImageRight{images/day11/PXL_20240819_151811404.MP.jpg}
%

%single
\landscapeThinMarginCentered{images/day11/PXL_20240819_152105639.jpg}
%

\pagestyle{fancy} % enable numbers